% Butterfly reduction with __shfl_xor_sync over eight lanes: three rounds
% with lane masks 4, 2 and 1; after the last round every lane holds the
% sum.  Original (the values are the chapter's own example).
% Usage: % Butterfly reduction with __shfl_xor_sync over eight lanes: three rounds
% with lane masks 4, 2 and 1; after the last round every lane holds the
% sum.  Original (the values are the chapter's own example).
% Usage: % Butterfly reduction with __shfl_xor_sync over eight lanes: three rounds
% with lane masks 4, 2 and 1; after the last round every lane holds the
% sum.  Original (the values are the chapter's own example).
% Usage: % Butterfly reduction with __shfl_xor_sync over eight lanes: three rounds
% with lane masks 4, 2 and 1; after the last round every lane holds the
% sum.  Original (the values are the chapter's own example).
% Usage: \input{figures/l05-butterfly}
% Width: body (120 mm); the drawing is about 104 mm wide.
\begin{tikzpicture}[x=1mm, y=1mm, font=\scriptsize\sffamily,
                    >={Stealth[length=1.2mm]}]
  \def\cw{11}
  \def\rh{11}
  \foreach \i in {0,...,7} {
    \node[text=ln-muted] at (\i*\cw+\cw/2,3*\rh+5) {\iflangja{レーン}{lane} \i};
  }
  % rows: round 0 (input), after xor 4, after xor 2, after xor 1
  \foreach \r/\lbl/\vals in {
      3/{\iflangja{入力}{input}}/{3,1,4,1,5,9,2,6},
      2/{$\oplus 4$}/{8,10,6,7,8,10,6,7},
      1/{$\oplus 2$}/{14,17,14,17,14,17,14,17},
      0/{$\oplus 1$}/{31,31,31,31,31,31,31,31}} {
    \node[anchor=east, text=ln-muted] at (-1,\r*\rh+2.5) {\lbl};
    \foreach \v [count=\i from 0] in \vals {
      \draw[fill=ln-box] (\i*\cw+1.5,\r*\rh) rectangle (\i*\cw+\cw-1.5,\r*\rh+5);
      \node at (\i*\cw+\cw/2,\r*\rh+2.5) {\v};
    }
  }
  % arrows: xor 4 (row 3 -> 2)
  \foreach \i in {0,1,2,3} {
    \draw[->, ln-accent] (\i*\cw+\cw/2,3*\rh-0.3) -- (\i*\cw+4*\cw+\cw/2,2*\rh+5.3);
    \draw[->, ln-accent] (\i*\cw+4*\cw+\cw/2,3*\rh-0.3) -- (\i*\cw+\cw/2,2*\rh+5.3);
  }
  % xor 2 (row 2 -> 1)
  \foreach \i in {0,1,4,5} {
    \draw[->, ln-accent] (\i*\cw+\cw/2,2*\rh-0.3) -- (\i*\cw+2*\cw+\cw/2,\rh+5.3);
    \draw[->, ln-accent] (\i*\cw+2*\cw+\cw/2,2*\rh-0.3) -- (\i*\cw+\cw/2,\rh+5.3);
  }
  % xor 1 (row 1 -> 0)
  \foreach \i in {0,2,4,6} {
    \draw[->, ln-accent] (\i*\cw+\cw/2,\rh-0.3) -- (\i*\cw+\cw+\cw/2,5.3);
    \draw[->, ln-accent] (\i*\cw+\cw+\cw/2,\rh-0.3) -- (\i*\cw+\cw/2,5.3);
  }
  \node[anchor=west, text=ln-muted] at (0,-4)
    {\iflangja{各段で，レーン $i$ は自分の値にレーン $i \oplus m$ の値を加える}{in each round lane $i$ adds the value of lane $i \oplus m$ to its own}};
\end{tikzpicture}

% Width: body (120 mm); the drawing is about 104 mm wide.
\begin{tikzpicture}[x=1mm, y=1mm, font=\scriptsize\sffamily,
                    >={Stealth[length=1.2mm]}]
  \def\cw{11}
  \def\rh{11}
  \foreach \i in {0,...,7} {
    \node[text=ln-muted] at (\i*\cw+\cw/2,3*\rh+5) {\iflangja{レーン}{lane} \i};
  }
  % rows: round 0 (input), after xor 4, after xor 2, after xor 1
  \foreach \r/\lbl/\vals in {
      3/{\iflangja{入力}{input}}/{3,1,4,1,5,9,2,6},
      2/{$\oplus 4$}/{8,10,6,7,8,10,6,7},
      1/{$\oplus 2$}/{14,17,14,17,14,17,14,17},
      0/{$\oplus 1$}/{31,31,31,31,31,31,31,31}} {
    \node[anchor=east, text=ln-muted] at (-1,\r*\rh+2.5) {\lbl};
    \foreach \v [count=\i from 0] in \vals {
      \draw[fill=ln-box] (\i*\cw+1.5,\r*\rh) rectangle (\i*\cw+\cw-1.5,\r*\rh+5);
      \node at (\i*\cw+\cw/2,\r*\rh+2.5) {\v};
    }
  }
  % arrows: xor 4 (row 3 -> 2)
  \foreach \i in {0,1,2,3} {
    \draw[->, ln-accent] (\i*\cw+\cw/2,3*\rh-0.3) -- (\i*\cw+4*\cw+\cw/2,2*\rh+5.3);
    \draw[->, ln-accent] (\i*\cw+4*\cw+\cw/2,3*\rh-0.3) -- (\i*\cw+\cw/2,2*\rh+5.3);
  }
  % xor 2 (row 2 -> 1)
  \foreach \i in {0,1,4,5} {
    \draw[->, ln-accent] (\i*\cw+\cw/2,2*\rh-0.3) -- (\i*\cw+2*\cw+\cw/2,\rh+5.3);
    \draw[->, ln-accent] (\i*\cw+2*\cw+\cw/2,2*\rh-0.3) -- (\i*\cw+\cw/2,\rh+5.3);
  }
  % xor 1 (row 1 -> 0)
  \foreach \i in {0,2,4,6} {
    \draw[->, ln-accent] (\i*\cw+\cw/2,\rh-0.3) -- (\i*\cw+\cw+\cw/2,5.3);
    \draw[->, ln-accent] (\i*\cw+\cw+\cw/2,\rh-0.3) -- (\i*\cw+\cw/2,5.3);
  }
  \node[anchor=west, text=ln-muted] at (0,-4)
    {\iflangja{各段で，レーン $i$ は自分の値にレーン $i \oplus m$ の値を加える}{in each round lane $i$ adds the value of lane $i \oplus m$ to its own}};
\end{tikzpicture}

% Width: body (120 mm); the drawing is about 104 mm wide.
\begin{tikzpicture}[x=1mm, y=1mm, font=\scriptsize\sffamily,
                    >={Stealth[length=1.2mm]}]
  \def\cw{11}
  \def\rh{11}
  \foreach \i in {0,...,7} {
    \node[text=ln-muted] at (\i*\cw+\cw/2,3*\rh+5) {\iflangja{レーン}{lane} \i};
  }
  % rows: round 0 (input), after xor 4, after xor 2, after xor 1
  \foreach \r/\lbl/\vals in {
      3/{\iflangja{入力}{input}}/{3,1,4,1,5,9,2,6},
      2/{$\oplus 4$}/{8,10,6,7,8,10,6,7},
      1/{$\oplus 2$}/{14,17,14,17,14,17,14,17},
      0/{$\oplus 1$}/{31,31,31,31,31,31,31,31}} {
    \node[anchor=east, text=ln-muted] at (-1,\r*\rh+2.5) {\lbl};
    \foreach \v [count=\i from 0] in \vals {
      \draw[fill=ln-box] (\i*\cw+1.5,\r*\rh) rectangle (\i*\cw+\cw-1.5,\r*\rh+5);
      \node at (\i*\cw+\cw/2,\r*\rh+2.5) {\v};
    }
  }
  % arrows: xor 4 (row 3 -> 2)
  \foreach \i in {0,1,2,3} {
    \draw[->, ln-accent] (\i*\cw+\cw/2,3*\rh-0.3) -- (\i*\cw+4*\cw+\cw/2,2*\rh+5.3);
    \draw[->, ln-accent] (\i*\cw+4*\cw+\cw/2,3*\rh-0.3) -- (\i*\cw+\cw/2,2*\rh+5.3);
  }
  % xor 2 (row 2 -> 1)
  \foreach \i in {0,1,4,5} {
    \draw[->, ln-accent] (\i*\cw+\cw/2,2*\rh-0.3) -- (\i*\cw+2*\cw+\cw/2,\rh+5.3);
    \draw[->, ln-accent] (\i*\cw+2*\cw+\cw/2,2*\rh-0.3) -- (\i*\cw+\cw/2,\rh+5.3);
  }
  % xor 1 (row 1 -> 0)
  \foreach \i in {0,2,4,6} {
    \draw[->, ln-accent] (\i*\cw+\cw/2,\rh-0.3) -- (\i*\cw+\cw+\cw/2,5.3);
    \draw[->, ln-accent] (\i*\cw+\cw+\cw/2,\rh-0.3) -- (\i*\cw+\cw/2,5.3);
  }
  \node[anchor=west, text=ln-muted] at (0,-4)
    {\iflangja{各段で，レーン $i$ は自分の値にレーン $i \oplus m$ の値を加える}{in each round lane $i$ adds the value of lane $i \oplus m$ to its own}};
\end{tikzpicture}

% Width: body (120 mm); the drawing is about 104 mm wide.
\begin{tikzpicture}[x=1mm, y=1mm, font=\scriptsize\sffamily,
                    >={Stealth[length=1.2mm]}]
  \def\cw{11}
  \def\rh{11}
  \foreach \i in {0,...,7} {
    \node[text=ln-muted] at (\i*\cw+\cw/2,3*\rh+5) {\iflangja{レーン}{lane} \i};
  }
  % rows: round 0 (input), after xor 4, after xor 2, after xor 1
  \foreach \r/\lbl/\vals in {
      3/{\iflangja{入力}{input}}/{3,1,4,1,5,9,2,6},
      2/{$\oplus 4$}/{8,10,6,7,8,10,6,7},
      1/{$\oplus 2$}/{14,17,14,17,14,17,14,17},
      0/{$\oplus 1$}/{31,31,31,31,31,31,31,31}} {
    \node[anchor=east, text=ln-muted] at (-1,\r*\rh+2.5) {\lbl};
    \foreach \v [count=\i from 0] in \vals {
      \draw[fill=ln-box] (\i*\cw+1.5,\r*\rh) rectangle (\i*\cw+\cw-1.5,\r*\rh+5);
      \node at (\i*\cw+\cw/2,\r*\rh+2.5) {\v};
    }
  }
  % arrows: xor 4 (row 3 -> 2)
  \foreach \i in {0,1,2,3} {
    \draw[->, ln-accent] (\i*\cw+\cw/2,3*\rh-0.3) -- (\i*\cw+4*\cw+\cw/2,2*\rh+5.3);
    \draw[->, ln-accent] (\i*\cw+4*\cw+\cw/2,3*\rh-0.3) -- (\i*\cw+\cw/2,2*\rh+5.3);
  }
  % xor 2 (row 2 -> 1)
  \foreach \i in {0,1,4,5} {
    \draw[->, ln-accent] (\i*\cw+\cw/2,2*\rh-0.3) -- (\i*\cw+2*\cw+\cw/2,\rh+5.3);
    \draw[->, ln-accent] (\i*\cw+2*\cw+\cw/2,2*\rh-0.3) -- (\i*\cw+\cw/2,\rh+5.3);
  }
  % xor 1 (row 1 -> 0)
  \foreach \i in {0,2,4,6} {
    \draw[->, ln-accent] (\i*\cw+\cw/2,\rh-0.3) -- (\i*\cw+\cw+\cw/2,5.3);
    \draw[->, ln-accent] (\i*\cw+\cw+\cw/2,\rh-0.3) -- (\i*\cw+\cw/2,5.3);
  }
  \node[anchor=west, text=ln-muted] at (0,-4)
    {\iflangja{各段で，レーン $i$ は自分の値にレーン $i \oplus m$ の値を加える}{in each round lane $i$ adds the value of lane $i \oplus m$ to its own}};
\end{tikzpicture}
